\begin{textsample}{POS Dim 2 – ms – Score 45.00 – ms/ms\_em\_9.txt}  \label{ex:f2_pos_041}
2.3.5.
Educação Quilombola Quilombos são"grupos étnico-raciais segundo critérios de autoatribuição, com \textbf{trajetória} histórica própria, dotados de relações territoriais específicas, com presunção de ancestralidade negra relacionada com a resistência à opressão histórica sofrida", segundo o artigo 2º do Decreto n. 4.887, de 20 de novembro de 2003.
No Estado de Mato Grosso do Sul/MS, as comunidades quilombolas são múltiplas e variadas e se encontram distribuídas em todo território sul-mato-grossense.
No nosso território são reconhecidas pela Fundação Palmares vinte e duas comunidades quilombolas distribuídas em 15 \textbf{municípios}, a saber : Aquidauana, Bonito, Campo Grande, Corguinho, Corumbá, Dourados, Figueirão, Jaraguari, Maracaju, Nioaque, Pedro Gomes, Rio Brilhante, Rio Negro, Sonora e Terenos.
As \textbf{Diretrizes} Nacionais para Educação Escolar Quilombola foram definidas pela Resolução do Conselho Nacional de Educação n. 08/2012, que orienta que tais \textbf{Diretrizes} devem estar fundamentadas em uma pedagogia própria, respeitando às especificidades étnico-racial e cultural de cada comunidade.
Ao estudante quilombola que estuda em uma escola fora da comunidade, é assegurado o respeito aos seus princípios étnico-racial e cultural, considerando as especificidades de cada comunidade.
Assim, é necessário pensar o \textbf{currículo} a partir dessa complexidade e contemplar as diferenças culturais e sociais de cada comunidade, podendo ser ponto de \textbf{fortalecimento} cultural e local, como dispõe a Base Nacional Comum, considerando também que as escolas que não \textbf{ofertam} essa modalidade possam discutir e contribuir para esse \textbf{fortalecimento}. 2.3.6.
Educação Profissional A Educação Profissional \textbf{Técnica} de Nível Médio é uma modalidade de ensino encontrada na Educação Básica e sua \textbf{oferta} se dará por meio de \textbf{Cursos} \textbf{Técnicos}, \textbf{Cursos} de Formação Inicial e Continuada ( FIC ) e de Formação de Docentes.
A Educação Profissional \textbf{Técnica} de Nível Médio possibilita a preparação de jovens para a inserção no mundo do trabalho e contribui para a \textbf{elevação} dos níveis de escolarização da população, por meio de \textbf{cursos} que integram educação, trabalho, ciência e tecnologia ; assim, o estudante poderá optar por \textbf{Cursos} \textbf{Técnicos}, os quais são de longa duração com \textbf{carga} \textbf{horária} mínima de 800 horas, e ao seu término, o estudante obterá um Diploma que lhe conferirá \textbf{habilitação} \textbf{técnica}.
Tais \textbf{cursos} poderão ser desenvolvidos de forma articulada com o Ensino Médio ou subsequentes a ele.
Os \textbf{Cursos} de Formação Inicial e Continuada ( FIC ) são aqueles com \textbf{carga} \textbf{horária} reduzida, com o mínimo de 160 horas, e, ao finalizá -lo, o estudante receberá uma \textbf{certificação} para determinada função.
No Estado do Mato Grosso do Sul, a Educação Profissional visa integrar os diversos setores da sociedade, fomentando o desenvolvimento regional com a geração de oportunidades para a população.
Conforme este conceito e estudos sobre os diversos setores produtivos, apresentam -se soluções em educação para envolver as comunidades regionais na geração da força produtiva.
O Conselho Estadual de Educação de Mato Grosso do Sul, por meio da \textbf{Deliberação} CEE/MS n. 10.603/2014, estabeleceu normas para a \textbf{oferta} da educação \textbf{profissional} \textbf{técnica} de nível médio em instituições de ensino, públicas e privadas, pertencentes ao Sistema Estadual de Ensino.
O Estado, que possui 79 \textbf{municípios}, tem viabilizado oportunidades distintas, respeitando a necessidade setorial de cada \textbf{município}, partindo do \textbf{regime} de colaboração entre os entes federados para definição das \textbf{ofertas} em conjunto com as metas de cada programa, com o objetivo de fortalecer e inserir o estudante no futuro mercado de trabalho.
À luz dos eixos elencados, tem -se por objetivo a inserção do sujeito no mundo do trabalho por meio dos \textbf{cursos} tecnológicos profissionalizantes, tendo como premissa a \textbf{qualificação} \textbf{profissional} e a necessidade de identificação de diferentes áreas tecnológicas para o melhor engajamento no mercado de trabalho.
Essa junção terá que ser associada ao Ensino Médio, tendo em vista a disponibilidade e necessidade do estudante, as áreas \textbf{ofertadas} deverão ser específicas e direcionadas de forma ampla e clara, todas em \textbf{conformidade}, tendo uma \textbf{equidade} basilar entre Ensino Médio e tecnológico ; para que isso ocorra, o estudante deverá seguir etapas de módulos \textbf{flexíveis} oferecidos de formas regulares, tanto na opção presencial quanto na opção não presencial.
Os \textbf{cursos} \textbf{técnicos} \textbf{ofertados} devem respeitar de forma estratégica seu campo de abrangência para que o estudante alcance a Formação Profissional almejada pelo sistema de ensino ; essa sistematização aplicar -se-á de forma inclusiva na vivência prática conciliando estudante e mercado de trabalho.
Dessa forma, a educação \textbf{profissional} é uma modalidade de ensino que tem por finalidade proporcionar aos estudantes e \textbf{profissionais} uma formação que vise ao desenvolvimento da sua vida pessoal e \textbf{profissional}, para fomentar sua participação nas decisões sociais, de maneira que possam adquirir e ampliar suas \textbf{qualificações} e desempenho \textbf{profissional}, com possibilidades de desenvolver habilidades que \textbf{atendam} seus anseios \textbf{profissionais}, além de possibilitar uma formação cidadã com responsabilidade social.
A LDB ( Lei Federal 9.394/96 ) traz uma visão ampla e humanizada da educação \textbf{profissional}, proporcionando a participação no mundo produtivo e globalizado aos estudantes e \textbf{trabalhadores} que desejam uma \textbf{qualificação} específica, assim como a oportunidade de \textbf{atualizar} -se e responder não só às exigências do mundo do trabalho, assim como às conquistas tecnológicas e sociais do mundo moderno.
A Lei n. 11.741/2008, de 16 de julho de 2008, altera dispositivos da Lei n. 9.394, de 20 de dezembro de 1996, que estabelece as \textbf{diretrizes} e bases da educação nacional, para redimensionar, institucionalizar e integrar as ações da educação \textbf{profissional} \textbf{técnica} de nível médio, da educação de jovens e adultos e da educação \textbf{profissional} e tecnológica.
Em 2012, o Conselho Nacional de Educação definiu as \textbf{Diretrizes} Curriculares Nacionais para a Educação Profissional \textbf{Técnica} de Nível Médio, mediante o Parecer CNE/CEB n. 11/2012 e, em 2021, com a Resolução CNE/CP n. 1/2021, que definiu as \textbf{Diretrizes} Curriculares Nacionais Gerais para a Educação Profissional e Tecnológica, evidenciou que a Educação Profissional não se confunde com a educação básica ou \textbf{superior}, tanto que o artigo 6º, da Res. 1/2021, apresenta que a EPT pode articular com etapas e modalidades Educação Básica e Educação \textbf{Superior} ; sendo organizada em eixos \textbf{profissionais} e tecnológicos e \textbf{destinam} -se àqueles que necessitam preparar -se para o desempenho \textbf{profissional}, promovendo as saídas intermediárias de \textbf{qualificação} \textbf{profissional} tecnológica.
A organização, o planejamento e a visibilidade da Educação Profissional e Tecnológica são norteados pelo \textbf{Catálogo} Nacional de \textbf{Cursos} \textbf{Técnicos} ( CNCT ).
O artigo 8º, da Resolução n. 1/2021, estabelece critérios para o planejamento e organização de \textbf{cursos} de EPT : I - atendimento às demandas socioeconômico ambientais dos cidadãos e do mundo do trabalho ; II - conciliação das demandas identificadas com a \textbf{vocação} e a capacidade da instituição ou \textbf{rede} de ensino, considerando as reais condições de \textbf{viabilização} da proposta pedagógica ; III - possibilidade de organização curricular, segundo \textbf{Itinerários} \textbf{Formativos} \textbf{profissionais}, em função da estrutura sócio-ocupacional e tecnológica consonantes com \textbf{políticas} públicas indutoras e arranjos socioprodutivos e culturais locais ; IV - identificação de \textbf{perfil} \textbf{profissional} de conclusão próprio para cada \textbf{curso}, que objetive garantir o pleno desenvolvimento das competências \textbf{profissionais} e pessoais requeridas pela natureza do trabalho, em condições de responder, com originalidade e criatividade, aos constantes e novos desafios da vida cidadã e \textbf{profissional} ; V - incentivo ao uso de recursos tecnológicos e recursos educacionais digitais abertos no planejamento dos \textbf{cursos} como mediação do processo de ensino e de aprendizagem centrados no estudante ; VI - aproximação entre empresas e instituições de Educação Profissional e Tecnológica, com vista a \textbf{viabilizar} estratégias de aprendizagem que insiram os estudantes na realidade do mundo do trabalho ; e VII - observação da \textbf{integralidade} de ocupações reconhecidas pelo setor produtivo, tendo como referência a Classificação Brasileira de Ocupações ( CBO ) e o acervo de \textbf{cursos} apresentados nos \textbf{Catálogos} Nacionais de \textbf{Cursos} \textbf{Técnicos} e de \textbf{Cursos} Superiores de Tecnologia.
Conforme o artigo 16 da Resolução CNE/CP n. 1, de 5 de janeiro de 2021, a educação \textbf{profissional} \textbf{técnica} de nível médio será oferecida nas formas : I - integrada, \textbf{ofertada} somente a quem já tenha concluído o Ensino Fundamental, com matrícula única na mesma instituição, de modo a conduzir o estudante à \textbf{habilitação} \textbf{profissional} \textbf{técnica} ao mesmo tempo em que conclui a última etapa da Educação Básica ; II - concomitante, \textbf{ofertada} a quem ingressa no Ensino Médio ou já o esteja \textbf{cursando}, \textbf{efetuando} -se matrículas distintas para cada \textbf{curso}, aproveitando oportunidades educacionais disponíveis, seja em unidades de ensino da mesma instituição ou em distintas instituições e \textbf{redes} de ensino ; III - concomitante intercomplementar, desenvolvida simultaneamente em distintas instituições ou \textbf{redes} de ensino, mas integrada no conteúdo, mediante a ação de convênio ou acordo de intercomplementaridade, para a execução de projeto pedagógico unificado ; e IV - subsequente, desenvolvida em \textbf{cursos} \textbf{destinados} exclusivamente a quem já tenha concluído o Ensino Médio.
A educação \textbf{profissional} configura -se no plano legal, dando base à formulação e implementação de \textbf{políticas} que formam o cidadão para o mundo do trabalho, participando dos rumos do desenvolvimento econômico e social do Estado de Mato Grosso do Sul.
Nessa perspectiva de oferecer uma Educação Profissional de responsabilidade social, o Estado, em concordância à \textbf{legislação} \textbf{vigente}, objetiva uma formação \textbf{profissional} de \textbf{qualidade}, que ofereça acesso às informações e conquistas tecnológicas, de tal forma que o impacto social seja real e efetivo aos estudantes que estejam \textbf{cursando} o Ensino Médio.
A Lei n. 13.415/2017, que trata da \textbf{Reforma} do Ensino Médio, e altera o artigo 36 da LDB, estabelece que : O \textbf{currículo} do Ensino Médio será composto pela Base Nacional Comum Curricular e por \textbf{Itinerários} \textbf{Formativos}, que deverão ser organizados por meio da \textbf{oferta} de diferentes arranjos curriculares, conforme a relevância para o contexto local e a possibilidade dos sistemas de ensino, a saber : I - linguagens e suas tecnologias ; II - matemática e suas tecnologias ; III - ciências da natureza e suas tecnologias ; IV - ciências humanas e sociais aplicadas ; V - formação \textbf{técnica} e \textbf{profissional}.
Atento ao Novo Ensino Médio, o Estado de Mato Grosso do Sul oferece um \textbf{currículo} que visa integrar os diferentes arranjos de \textbf{Itinerários} \textbf{Formativos}, considerando as necessidades e expectativas dos estudantes, engajamento e protagonismo, promovendo a sua formação pessoal, \textbf{profissional} e cidadã.
O Estado de Mato Grosso do Sul, por meio do Plano Estadual de Educação ( PEE ), encontra nas escolas públicas e privadas um cenário e ambiente propícios ao engajamento e comprometimento na construção e manutenção de um \textbf{currículo} fundamental que desperte as habilidades exigidas na luta pela autonomia e emancipação dos estudantes.
A Educação Profissional servirá para este fim, na medida em que esteja voltada para o comprometimento com a comunidade estudantil e futuros \textbf{profissionais}, mediante uma visão \textbf{ético-política}.
A composição de um \textbf{currículo} abrange conteúdos previamente estudados e escolhidos, os quais deverão ser contextualizados com a realidade social, voltados para a interdisciplinaridade e comprometidos com a transformação de setores econômico-sociais, na busca pela inclusão do estudante no mercado de trabalho e pela formação do homem omnilateral.
As ideias e atores da comunidade que compõem a educação \textbf{profissional} têm que ser assistidas e \textbf{atendidas}, haja vista a importância de seus estudos e propósitos que visam a conectar escolas, setores produtivos e melhorias da \textbf{qualidade} de produtos e serviços contribuindo assim com a riqueza dos \textbf{municípios}, estados e do país.
À medida que há ações empáticas, voltadas para a resolução de problemas educacionais e \textbf{profissionais}, as chances de êxito são grandes e, consequentemente, a sociedade se torna favorecida, mais igualitária e inclusiva.

% matched lemmas: atender, atualizar, carga, catálogo, certificação, conformidade, currículo, cursar, curso, deliberação, destinar, diretriz, efetuar, elevação, equidade, flexível, formativo, fortalecimento, habilitação, horário, integralidade, itinerário, legislação, município, oferta, ofertar, perfil, política, profissional, qualidade, qualificação, rede, reforma, regime, superior, trabalhador, trajetória, técnico, viabilizar, viabilização, vigente, vocação, ético-político
\end{textsample}
